\documentclass[12pt]{ximera}
\title{Exam 1}
\begin{document}
\begin{abstract}
Exam 1, covering Chapter 1: majority and Condorcet candidates, Independence of Irrelevant
Alternatives, a truncated Plurality-with-Elimination election, and bookkeeping for
Pairwise Comparisons points.
\end{abstract}
\maketitle

\section*{Exam 1}

\begin{problem}
Assess whether each of the following claims is true or false.

\begin{prompt}
\textbf{(a)} Claim: if a candidate has a majority of first-place votes, then they are a
Condorcet candidate.

\begin{multipleChoice}
\choice[correct]{True}
\choice{False}
\end{multipleChoice}

\begin{feedback}
True. In any head-to-head matchup, every voter who ranks the candidate first prefers
them to the opponent. A majority of first-place votes is therefore a majority in every
matchup.
\end{feedback}
\end{prompt}

\begin{prompt}
\textbf{(b)} Claim: a candidate with a majority of first-place votes will always win an
election using the Pairwise Comparisons method.

\begin{multipleChoice}
\choice[correct]{True}
\choice{False}
\end{multipleChoice}

\begin{feedback}
True. By part (a) such a candidate is a Condorcet candidate, so they win every
comparison and finish with the most points.
\end{feedback}
\end{prompt}
\end{problem}

\begin{problem}
Consider the following preference schedule for an election with four candidates.

\begin{center}
\begin{tabular}{c|ccc}
Preference & 10 & 8 & 6 \\\hline
1st & C & A & B \\
2nd & D & B & D \\
3rd & A & C & C \\
4th & B & D & A
\end{tabular}
\end{center}

\begin{prompt}
\textbf{(a)} Give the complete ranking by the Plurality method.

1st: $\answer{C}$ with $\answer{10}$ \quad 2nd: $\answer{A}$ with $\answer{8}$ \quad
3rd: $\answer{B}$ with $\answer{6}$ \quad 4th: $\answer{D}$ with $\answer{0}$

\begin{feedback}
First-place votes: C $10$, A $8$, B $6$, D $0$.
\end{feedback}
\end{prompt}

\begin{prompt}
\textbf{(b)} Which candidate, if any, can be removed to illustrate that the Plurality
method violates the Independence-of-Alternatives criterion?

\begin{multipleChoice}
\choice[correct]{A}
\choice{B}
\choice{C}
\choice{D}
\choice{None of these}
\end{multipleChoice}

\begin{feedback}
Remove \textbf{A}. The $8$ voters who ranked A first now rank B first, so the counts
become B $6+8=14$ and C $10$ --- and B wins instead of C. Removing B or D leaves C
winning, and C is the winner, so C cannot be the one removed.
\end{feedback}
\end{prompt}

\begin{prompt}
\textbf{(c)} Is there a Condorcet candidate in this election?

\begin{multipleChoice}
\choice{Yes, A}
\choice{Yes, B}
\choice{Yes, C}
\choice{Yes, D}
\choice[correct]{No}
\end{multipleChoice}

\begin{feedback}
No. The head-to-head results are
\[
\begin{array}{lll}
\text{A vs B} & 18\text{--}6 & \text{A}\\
\text{A vs C} & 8\text{--}16 & \text{C}\\
\text{A vs D} & 8\text{--}16 & \text{D}\\
\text{B vs C} & 14\text{--}10 & \text{B}\\
\text{B vs D} & 14\text{--}10 & \text{B}\\
\text{C vs D} & 18\text{--}6 & \text{C}
\end{array}
\]
Every candidate loses at least once: A beats B, B beats C, and C beats A --- a cycle.
So no one beats everyone.
\end{feedback}
\end{prompt}
\end{problem}

\begin{problem}
Consider the following truncated preference schedule for an election with four
candidates.

\begin{center}
\begin{tabular}{c|ccccc}
Preference & 9 & 7 & 5 & 3 & 1 \\\hline
1st & B & D & C & A & C \\
2nd & C & A & D & D & B
\end{tabular}
\end{center}

\begin{prompt}
\textbf{(a)} How many voters participated? $\answer{25}$

\begin{feedback}
$9+7+5+3+1=25$.
\end{feedback}
\end{prompt}

\begin{prompt}
\textbf{(b)} Using Plurality-with-Elimination, find the Round 1 totals and the candidate
eliminated.

A: $\answer{3}$ \quad B: $\answer{9}$ \quad C: $\answer{6}$ \quad D: $\answer{7}$ \qquad
Eliminated: $\answer{A}$

\begin{feedback}
First-place votes: B $9$, D $7$, C $5+1=6$, A $3$. A has the fewest.
\end{feedback}
\end{prompt}

\begin{prompt}
\textbf{(c)} Round 2 totals and the candidate eliminated.

B: $\answer{9}$ \quad C: $\answer{6}$ \quad D: $\answer{10}$ \qquad Eliminated: $\answer{C}$

\begin{feedback}
A's $3$ voters ranked D second, so D rises to $10$. C now has the fewest.
\end{feedback}
\end{prompt}

\begin{prompt}
\textbf{(d)} Finish the method and give the complete ranking.

Round 3 --- B: $\answer{10}$ \quad D: $\answer{15}$

\smallskip
1st $\answer{D}$, 2nd $\answer{B}$, 3rd $\answer{C}$, 4th $\answer{A}$

\begin{feedback}
C's votes transfer: the $5$ voters (C, D) go to D and the $1$ voter (C, B) goes to B,
giving D $15$ and B $10$. D wins; the ranking is D, B, C, A.
\end{feedback}
\end{prompt}

\begin{prompt}
\textbf{(e)} How many ballots were exhausted? $\answer{0}$

\begin{feedback}
None. Every eliminated candidate's voters had a surviving second choice when their
ballot was transferred.
\end{feedback}
\end{prompt}
\end{problem}

\begin{problem}
Suppose an election with 6 candidates ($A,B,C,D,E,F$) is run using the Pairwise Comparisons
method.

\begin{prompt}
\textbf{(a)} How many pairwise comparisons will be necessary?

$\answer{15}$

\begin{feedback}
Each pair of candidates meets once: $C(6,2)=\frac{6\cdot 5}{2}=15$.
\end{feedback}
\end{prompt}

\begin{prompt}
\textbf{(b)} What is the total number of Pairwise Comparisons points awarded? (Each
comparison awards $1$ point to the winner, or half a point to each of two tied
candidates.)

$\answer{15}$

\begin{feedback}
Every comparison hands out exactly $1$ point in total, whether it is won or tied, so
$15$ comparisons award $15$ points.
\end{feedback}
\end{prompt}

\begin{prompt}
\textbf{(c)} What is the maximum number of points a single candidate can earn?

$\answer{5}$

\begin{feedback}
A candidate takes part in only $5$ comparisons --- one against each other candidate
--- so the most they can earn is $5$, by winning them all.
\end{feedback}
\end{prompt}

\begin{prompt}
\textbf{(d)} Suppose candidate A wins $3$ comparisons and ties none, B wins $1$ and ties $1$, C wins $1$ and
ties $2$, D wins $0$ and ties $2$, and E is a Condorcet candidate. How many points does F earn?

$\answer{2.5}$

\begin{hint}
A Condorcet candidate wins every comparison. Then use the total from part (b).
\end{hint}

\begin{feedback}
A earns $3$, B earns $1.5$, C earns $1+1=2$, D earns $1$, and E, as a Condorcet
candidate, wins all $5$ of its comparisons for $5$ points. That is $3+1.5+2+1+5=12.5$. The points total $15$, so F earns what is left: $2.5$ points.
\end{feedback}
\end{prompt}

\begin{prompt}
\textbf{(e)} Give the complete ranking of the candidates.

1st: $\answer{E}$ \quad 2nd: $\answer{A}$ \quad 3rd: $\answer{F}$ \quad 4th: $\answer{C}$ \quad 5th: $\answer{B}$ \quad 6th: $\answer{D}$

\begin{feedback}
By points: E $5$, A $3$, F $2.5$, C $2$, B $1.5$, D $1$.
\end{feedback}
\end{prompt}
\end{problem}

\end{document}
